% ECONOMICS_PROBLEM: {"id": "F36-407", "title": "International risk sharing", "jel_code": "F36", "source_id": "OP-0335"}
% !TeX program = xelatex
\documentclass[11pt,a4paper]{article}
\usepackage[margin=23mm]{geometry}
\usepackage{fontspec}
\setmainfont{Latin Modern Roman}
\usepackage{amsmath,amssymb,mathtools}
\usepackage{xcolor,enumitem,booktabs,longtable,array,xurl,hyperref}
\definecolor{muted}{HTML}{53636C}
\hypersetup{colorlinks=true,urlcolor=blue,linkcolor=blue}
\setlength{\parindent}{0pt}
\setlength{\parskip}{5pt}
\newcommand{\R}{\mathbb R}
\newcommand{\E}{\mathbb E}
\newcommand{\N}{\mathbb N}
\newcommand{\ind}{\mathbf 1}
\newcommand{\Var}{\operatorname{Var}}
\newcommand{\Cov}{\operatorname{Cov}}
\newcommand{\supp}{\operatorname{supp}}
\DeclareMathOperator*{\argmin}{arg\,min}
\DeclareMathOperator*{\argmax}{arg\,max}
\newcommand{\entrymeta}[3]{{\sffamily\small\color{muted}\textbf{\texttt{#1}}\quad|\quad #2\par #3\par}}
\newcommand{\Problemref}[1]{\texttt{#1}}
\newcommand{\JELref}[2]{\texttt{#2} (JEL \texttt{#1})}
\begin{document}
\section*{International risk sharing}
\textbf{Problem ID:} \texttt{F36-407}\quad \textbf{JEL:} \texttt{F36}\par
% BEGIN ECONOMICS STATEMENT
\label{problem:OP-0335}
{\sffamily\small\color{muted}Original subject group: International finance and exchange rates\par}
\label{definitions:problem:30007617}
\textbf{Audit classification:} Explicit published open question. Footnote 20 explicitly leaves a US/advanced-economy/emerging-market extension to future work. The catalogue variance-insurance statistic is editorial and does not itself measure welfare. Add a genuinely source-motivated three-country variant.
\subsection*{Definitions and precise research target}
Fix countries with household consumption \(c_{ct}>0\), price levels, bilateral asset holdings and returns. Let \(g_{ct}=\log c_{ct}-\log c_{c,t-1}\), and define country-specific consumption innovation as the residual \(e_{ct}\) after projecting \(g_{ct}\) on a declared common-information set. An integration policy \(a\) changes feasible cross-border asset trade. Define its consumption-insurance effect as
\[\Delta_c(a)=\operatorname{Var}(e_{ct}(do(a)))-\operatorname{Var}(e_{ct}(do(0))).\]
The baseline \(0\), sample population and aggregate controls must be fixed in advance. The task is to identify this effect and characterize how safe-asset demand and asymmetric risk-bearing redistribute consumption losses in crises. Model exchange-rate valuation changes, nominal rigidities and incomplete markets explicitly before interpreting aggregate consumption comovement as insurance. Separate country output shocks from preference and common global shocks; consumption correlations alone need not measure risk sharing. Determine which exposure or policy changes identify the counterfactual and account for endogenous portfolio adjustments. Report distributional effects within countries where aggregate gains conceal household losses. Complete-market equal-marginal-utility conditions are useful benchmarks, not assumptions that may be imposed to answer why real-world insurance is uneven.

\textbf{Additional formal definitions and scope (editorial).} The innovation must be defined relative to a fixed predictor space. One choice is \(e_{ct}(a)=g_{ct}(a)-\operatorname{Proj}_{\mathcal H}g_{ct}(a)\) in \(L^2\), where \(\mathcal H\) is a declared common-information linear subspace including constants; a conditional-expectation residual is a different nonlinear convention. The projection operator is re-estimated under each counterfactual law if the target is unexplained country-specific variance; holding estimated coefficients fixed yields another estimand. Consumption is per capita real consumption with a fixed deflator and household coverage. Lower residual variance measures this variance-insurance statistic and need not imply higher welfare if mean consumption falls or common shocks grow. Specify consumption utility, asset feasibility and country/household welfare before claiming welfare risk sharing.

For the identification part, declare an admissible class \(\Theta\) of structural models, including preferences, technologies, shock laws, measurement/sampling rules and any equilibrium selector. If \(O\) denotes the complete observable history and \(T(\theta)\) the target defined above, the identified set is \(\mathcal I_T=\{T(\theta):\theta\in\Theta,\ \mathcal L_\theta(O)=\mathcal L_{\rm obs}(O)\}\); point identification means this set is a singleton. A \(do(a)\) comparison replaces stated policy/technology equations while fixing the declared exogenous law and re-solving the remaining equations. These definitions do not assert that an identification theorem holds without further restrictions.
\subsection*{Variants and assumption changes}
\begin{enumerate}
\item \textbf{Complete-market benchmark (editorial).} State common beliefs, tradeable contingent claims and household marginal utilities; derive the insurance benchmark under those assumptions.
\item \textbf{Three-country extension (explicitly source-motivated).} Distinguish the US, other advanced economies and emerging markets, with their monetary regimes and shock laws stated; study safety-shock transmission and international insurance. Compare variance statistics with separate household welfare effects.
\end{enumerate}
\subsection*{References and source audit}
\begin{enumerate}
\item Primary January 2024 cover verifies Kekre and Lenel and title; earlier WP 29238 versions used different titles. Locator: Section 4, footnote 20, printed pp.25--26 (PDF pp.26--27). Support: Explicit three-country extension. Full January 2024 manuscript accessible. URL: \url{https://www.aeaweb.org/conference/2024/program/paper/zak85S66}.
\item AEA verifies Kekre and Lenel, 2024, AER 114(6):1650--1691 and DOI \nolinkurl{10.1257/aer.20211319}. Locator: Publisher citation and abstract; manuscript-specific open passage is located in the preceding reference. Published full-text access was restricted; publisher metadata/abstract inspected. URL: \url{https://www.aeaweb.org/articles?id=10.1257/aer.20211319}.
\end{enumerate}
\begin{enumerate}\raggedright
\item\hypertarget{definitions:source:30007617:1}{} Rohan Kekre; Moritz Lenel. 2024. \emph{The Flight to Safety and International Risk Sharing}. Section 4, footnote 20, printed pp.25--26 (PDF pp.26--27).
\url{https://www.aeaweb.org/conference/2024/program/paper/zak85S66}.
\item\hypertarget{definitions:source:30007617:2}{} Rohan Kekre; Moritz Lenel. 2024. \emph{The Flight to Safety and International Risk Sharing}. Publisher citation and abstract; manuscript-specific open passage is located in the preceding reference.
\url{https://www.aeaweb.org/articles?id=10.1257/aer.20211319}.
DOI: \url{https://doi.org/10.1257/aer.20211319}.
\end{enumerate}

\subsection*{Common conventions: Source mathematical conventions}

These conventions apply to each entry unless explicitly replaced.

\textbf{Models and identification.} A primitive model \(m\) specifies all probability laws, feasible choices, information, equilibrium and measurement rules used by its target. The admissible class \(\mathcal M\) is fixed independently of outcomes. If \(P_O(m)\) is its observable probability law and \(\tau(m)\) the target, its identified set at observable law \(P\) is
\[
 \mathcal I_\tau(P)=\{\tau(m):m\in\mathcal M,\ P_O(m)=P\}.
\]
Point identification means that this set is a singleton. An empty set signals incompatibility of data and assumptions. Coordinate bounds need not describe the joint set. Bounds are sharp only if every retained target is attainable or arbitrarily approachable by compatible models. Finite bounds require explicit restrictions on unobserved quantities; unrestricted confounding, hidden wealth, extrapolation or arbitrary equilibrium selection can yield uninformative sets.

\textbf{Causal interventions.} A potential outcome \(Y_i(p)\) is the outcome under a complete policy regime \(p\), including timing, financing, information and market spillovers. The target population and units are fixed before treatment. With interference, use the complete assignment vector or a justified exposure mapping; individual no-interference notation alone cannot identify an economy-wide intervention. A difference of mean potential outcomes does not require the cross-world joint distribution. Conditioning on post-treatment migration, survival, enrollment or employment changes the estimand and needs separate justification.

\textbf{Statistical statements.} An estimator, confidence set, forecast or test specifies a sampling law, model class and loss/coverage criterion. Uniform \((1-\alpha)\)-coverage means \(\inf_{m\in\mathcal M}\Pr_m\{\tau(m)\in C_n\}\geq1-\alpha\), or an explicitly asymptotic version. The whole parameter space trivially has coverage, so informativeness requires a stated width, risk or power objective. A forecasting improvement is relative to a named benchmark, information set and evaluation population.

\textbf{Equilibrium and optimization.} An equilibrium comprises feasible plans, optimality, beliefs where needed, budget constraints and all market/resource clearing. The primitives or a stated benchmark define each correspondence \(\mathcal E(p;m)\). Nonemptiness is assumed or proved, never inferred just from compact policy sets. A selected-equilibrium target uses a declared selection rule; a robust target quantifies over all permitted equilibria. Use suprema or infima unless attainment is established. An \(\varepsilon\)-optimal policy is feasible and within \(\varepsilon\) of the finite optimum value. Report an empty feasible set separately.

\textbf{Welfare and accounting.} Normative utility, interpersonal normalization, social weights, numeraire, discounting and the use of public receipts are primitives. Behavioral utility need not coincide with the welfare criterion. Household welfare includes ownership income and rents once; adding profits again to compensating variation generally double-counts them. A monetary equivalent is defined by an explicit transfer and price convention, with monotonicity and range conditions. Average effects, mechanism contributions, transfers and welfare losses are different targets. Infinite sums require convergence and dynamic feasible plans require terminal or no-Ponzi conditions.

\textbf{Source versions.} A working-paper date, revision date and journal publication year are distinguished. A DOI for a journal version is not automatically the DOI of an earlier manuscript. Locators refer to the stated version and printed pages where specified. A metadata-only check does not verify a claim about a full-text conclusion. Every entry's source subsection narrows the provenance claim accordingly.
% END ECONOMICS STATEMENT
\end{document}
